\textbf{Phần 1. Trắc nghiệm lựa chọn}
\setcounterpageref{ex}{0}

\begin{ex} %[3D10N56]
Kết quả của phép tính $2\,000 \times 4$ là
\choice
{$6000$}
{\True $8000$}
{$10000$}
{$7000$}
\loigiai{
$2000 \times 4 = 8000$.\\
\textbf{Đáp án của câu hỏi là B}
}
\end{ex}

\begin{ex} %[3D10N56]
Kết quả của phép nhân không nhớ $1\,231 \times 3$ là
\choice
{$3683$}
{$3793$}
{\True $3693$}
{$3663$}
\loigiai{
$1231 \times 3 = 3693$.\\
\textbf{Đáp án của câu hỏi là C}
}
\end{ex}

\begin{ex} %[3D10H56]
Kết quả của phép nhân có nhớ $1\,418 \times 4$ là
\choice
{$5662$}
{$5682$}
{$5572$}
{\True $5672$}
\loigiai{
$1418 \times 4 = 5672$.\\
\textbf{Đáp án của câu hỏi là D}
}
\end{ex}

\begin{ex} %[3D10H56]
Kết quả của phép nhân $2\,315 \times 3$ là
\choice
{\True $6945$}
{$6935$}
{$6955$}
{$6845$}
\loigiai{
$2315 \times 3 = 6945$.\\
\textbf{Đáp án của câu hỏi là A}
}
\end{ex}

\begin{ex} %[3D10H56]
Tìm $x$ biết: $x : 3 = 1\,250$.
\choice
{$3650$}
{$3850$}
{$3550$}
{\True $3750$}
\loigiai{
$x = 1250 \times 3 = 3750$.\\
\textbf{Đáp án của câu hỏi là D}
}
\end{ex}

\begin{ex} %[3D10H56]
Tính giá trị biểu thức: $1\,200 \times 4 + 1\,500$.
\choice
{$6100$}
{\True $6300$}
{$6500$}
{$6000$}
\loigiai{
$4800 + 1500 = 6300$.\\
\textbf{Đáp án của câu hỏi là B}
}
\end{ex}

\begin{ex} %[3D10V56]
Một công ty may mỗi ngày may được $1\,150$ bộ quần áo. Hỏi trong 5 ngày công ty đó may được bao nhiêu bộ quần áo?
\choice
{$5650$ bộ}
{$5850$ bộ}
{\True $5750$ bộ}
{$5550$ bộ}
\loigiai{
$1150 \times 5 = 5750$ bộ.\\
\textbf{Đáp án của câu hỏi là C}
}
\end{ex}

\begin{ex} %[3D10V56]
Một cuộn dây thép dài $1\,420\text{ m}$. Hỏi 4 cuộn dây thép như thế dài tất cả bao nhiêu mét?
\choice
{$5580\text{ m}$}
{\True $5680\text{ m}$}
{$5780\text{ m}$}
{$5600\text{ m}$}
\loigiai{
$1420 \times 4 = 5680\text{ m}$.\\
\textbf{Đáp án của câu hỏi là B}
}
\end{ex}

\begin{ex} %[3D10V56]
Một hình vuông có cạnh dài $1\,250\text{ m}$. Chu vi hình vuông đó là
\choice
{\True $5000\text{ m}$ (hay $5\text{ km}$)}
{$4800\text{ m}$}
{$5200\text{ m}$}
{$4500\text{ m}$}
\loigiai{
$1250 \times 4 = 5000\text{ m}$.\\
\textbf{Đáp án của câu hỏi là A}
}
\end{ex}

\begin{ex} %[3D10V56]
Tính giá trị biểu thức: $2\,150 \times 4 - 3\,600$.
\choice
{$4800$}
{$5200$}
{\True $5000$}
{$5500$}
\loigiai{
$8600 - 3600 = 5000$.\\
\textbf{Đáp án của câu hỏi là C}
}
\end{ex}

\begin{ex} %[3D10V56]
Một xe tải chở $2\,350\text{ kg}$ hàng. Hỏi 3 xe tải như thế chở được bao nhiêu ki-lô-gam hàng?
\choice
{$6950\text{ kg}$}
{$7150\text{ kg}$}
{$6850\text{ kg}$}
{\True $7050\text{ kg}$}
\loigiai{
$2350 \times 3 = 7050\text{ kg}$.\\
\textbf{Đáp án của câu hỏi là D}
}
\end{ex}

\begin{ex} %[3D10C56]
Một cửa hàng nhập về 4 thùng bánh, mỗi thùng có $1\,200$ cái bánh. Cửa hàng đã bán được $3\,500$ cái. Số cái bánh còn lại là
\choice
{\True $1300$ cái}
{$1200$ cái}
{$1400$ cái}
{$1500$ cái}
\loigiai{
Tổng: $1200 \times 4 = 4800$ cái.\\
Còn lại: $4800 - 3500 = 1300$ cái bánh.\\
\textbf{Đáp án của câu hỏi là A}
}
\end{ex}

\textbf{Phần 2. Trắc nghiệm đúng sai}
\setcounterpageref{ex}{0}

\begin{ex} %[3D10H56]
Xét tính đúng sai của các phép nhân:
\choiceTF
{\True $1200 \times 4 = 4800$.}
{\True $2314 \times 2 = 4628$.}
{$1500 \times 5 = 7000$.}
{\True $1500 \times 5 = 7500$.}
\loigiai{
\begin{itemize}
    \item $1200 \times 4 = 4800$: Đúng.
    \item $2314 \times 2 = 4628$: Đúng.
    \item $1500 \times 5 = 7500 \neq 7000$. Ý 3 Sai, Ý 4 Đúng.
\end{itemize}
\begin{itemize}
    \item a) \textbf{Đúng}.
    \item b) \textbf{Đúng}.
    \item c) \textbf{Sai}.
    \item d) \textbf{Đúng}.
\end{itemize}
}
\end{ex}

\begin{ex} %[3D10V56]
Một trang trại nuôi 3 chuồng lợn, mỗi chuồng có $1\,120$ con. Xét tính đúng sai:
\choiceTF
{\True Cả 3 chuồng có tất cả $3\,360$ con lợn.}
{\True Nếu bán đi $1\,200$ con thì còn lại $2\,160$ con lợn.}
{Nếu nhập thêm 500 con thì tổng số lợn là $3\,800$ con.}
{\True Nếu nhập thêm 500 con thì tổng số lợn là $3\,860$ con ($3360 + 500 = 3860$).}
\loigiai{
\begin{itemize}
    \item Tổng: $1120 \times 3 = 3360$ con. Đúng.
    \item Còn lại: $3360 - 1200 = 2160$ con. Đúng.
    \item $3360 + 500 = 3860$ con $\neq 3800$. Ý 3 Sai, Ý 4 Đúng.
\end{itemize}
\begin{itemize}
    \item a) \textbf{Đúng}.
    \item b) \textbf{Đúng}.
    \item c) \textbf{Sai}.
    \item d) \textbf{Đúng}.
\end{itemize}
}
\end{ex}

\textbf{Phần 3. Trả lời ngắn}
\setcounterpageref{ex}{0}

\begin{ex} %[3D10H56]
Tính tích: $1\,312 \times 3$.
\shortans{3936}
\loigiai{
\textbf{Đáp số:} $3936$\\
$1312 \times 3 = 3936$.
}
\end{ex}

\begin{ex} %[3D10H56]
Tính tích: $2\,418 \times 4$.
\shortans{9672}
\loigiai{
\textbf{Đáp số:} $9672$\\
$2418 \times 4 = 9672$.
}
\end{ex}

\begin{ex} %[3D10V56]
Tính giá trị biểu thức: $1\,500 \times 4 - 2\,000$.
\shortans{4000}
\loigiai{
\textbf{Đáp số:} $4000$\\
$6000 - 2000 = 4000$.
}
\end{ex}

\begin{ex} %[3D10V56]
Một hình vuông có cạnh dài $1\,800\text{ m}$. Tính chu vi hình vuông theo mét.
\shortans{7200}
\loigiai{
\textbf{Đáp số:} $7200$\\
$1800 \times 4 = 7200\text{ m}$.
}
\end{ex}

\begin{ex} %[3D10V56]
Tìm $x$ biết: $x : 5 = 1\,420$.
\shortans{7100}
\loigiai{
\textbf{Đáp số:} $7100$\\
$x = 1420 \times 5 = 7100$.
}
\end{ex}

\begin{ex} %[3D10C56]
Một hình chữ nhật có chiều rộng $1\,150\text{ m}$, chiều dài gấp 3 lần chiều rộng. Tính chiều dài hình chữ nhật theo mét.
\shortans{3450}
\loigiai{
\textbf{Đáp số:} $3450$\\
$1150 \times 3 = 3450\text{ m}$.
}
\end{ex}
